\textbf{Phần 1. Trắc nghiệm lựa chọn}
\setcounterpageref{ex}{0}

\begin{ex} %[5D12V68]
Cho bốn chữ số $0, 2, 5, 8$. Viết được tất cả bao nhiêu số thập phân có đủ bốn chữ số trên mà phần thập phân có đúng một chữ số?
\choice
{$12\text{ số}$}
{\True $18\text{ số}$}
{$24\text{ số}$}
{$16\text{ số}$}
\loigiai{
Một số thập phân có $4$ chữ số khác nhau gồm $3$ chữ số ở phần nguyên và $1$ chữ số ở phần thập phân.\\
- Chọn chữ số hàng phần mười: có $4$ cách chọn.\\
- Chọn chữ số hàng trăm (khác $0$ và khác hàng phần mười): có $2$ cách chọn nếu hàng phần mười là $0$, hoặc $2$ cách nếu hàng phần mười khác $0$.
Cụ thể:
+ Trường hợp 1: Chữ số hàng phần mười là chữ số $0$ ($1$ cách). Chữ số hàng trăm có $3$ cách chọn, hàng chục có $2$ cách chọn, hàng đơn vị có $1$ cách chọn $\Rightarrow 1 \times 3 \times 2 \times 1 = 6$ số.
+ Trường hợp 2: Chữ số hàng phần mười khác $0$ ($3$ cách: $2, 5, 8$). Chữ số hàng trăm phải khác $0$ và khác hàng phần mười $\Rightarrow$ có $2$ cách chọn. Chữ số hàng chục có $2$ cách chọn, hàng đơn vị có $1$ cách chọn $\Rightarrow 3 \times 2 \times 2 \times 1 = 12$ số.
Tổng số các số viết được là:
$$6 + 12 = 18\text{ (số)}.$$
\textbf{Đáp án của câu hỏi là B}
}
\end{ex}

\begin{ex} %[5D12V68]
Phân số nào sau đây lớn hơn $\dfrac{3}{5}$ nhưng nhỏ hơn $\dfrac{4}{5}$?
\choice
{$\dfrac{1}{2}$}
{\True $\dfrac{7}{10}$}
{$\dfrac{4}{7}$}
{$\dfrac{5}{6}$}
\loigiai{
Quy đồng mẫu số chung là $10$:\\
$\dfrac{3}{5} = \dfrac{6}{10}$; $\dfrac{4}{5} = \dfrac{8}{10}$.\\
Phân số nằm giữa $\dfrac{6}{10}$ và $\dfrac{8}{10}$ là $\dfrac{7}{10}$.\\
\textbf{Đáp án của câu hỏi là B}
}
\end{ex}

\begin{ex} %[5D12V68]
Khi dịch dấu phẩy của một số thập phân sang bên phải một chữ số thì số đó tăng thêm $24{,}75$ đơn vị. Số thập phân ban đầu là
\choice
{$2{,}475$}
{\True $2{,}75$}
{$0{,}275$}
{$27{,}5$}
\loigiai{
Dịch dấu phẩy sang phải một chữ số tức là nhân số đó với $10$.\\
Số mới gấp $10$ lần số ban đầu, tức là tăng thêm $9$ lần số ban đầu.\\
Số thập phân ban đầu là:
$$24{,}75 : (10 - 1) = 24{,}75 : 9 = 2{,}75.$$
\textbf{Đáp án của câu hỏi là B}
}
\end{ex}

\begin{ex} %[5D12V68]
Tìm số tự nhiên $x$ biết: $3{,}75 < \dfrac{x}{4} < 4{,}25$.
\choice
{$14$}
{\True $16$}
{$15$}
{$18$}
\loigiai{
Nhân cả ba vế với $4$:
$$3{,}75 \times 4 < x < 4{,}25 \times 4$$
$$15 < x < 17$$
Vì $x$ là số tự nhiên nên $x = 16$.\\
\textbf{Đáp án của câu hỏi là B}
}
\end{ex}

\begin{ex} %[5D12V68]
Viết phân số $\dfrac{13}{25}$ dưới dạng số thập phân ta được
\choice
{$0{,}13$}
{\True $0{,}52$}
{$0{,}26$}
{$0{,}54$}
\loigiai{
Ta có: $\dfrac{13}{25} = \dfrac{13 \times 4}{25 \times 4} = \dfrac{52}{100} = 0{,}52$.\\
\textbf{Đáp án của câu hỏi là B}
}
\end{ex}

\begin{ex} %[5D12V68]
Sắp xếp các số sau theo thứ tự từ bé đến lớn: $3{,}14;\; 3\dfrac{1}{5};\; 3{,}125;\; \dfrac{13}{4}$.
\choice
{$3{,}125 < 3{,}14 < \dfrac{13}{4} < 3\dfrac{1}{5}$}
{\True $3{,}125 < 3{,}14 < 3\dfrac{1}{5} < \dfrac{13}{4}$}
{$3{,}14 < 3{,}125 < 3\dfrac{1}{5} < \dfrac{13}{4}$}
{$3{,}125 < 3\dfrac{1}{5} < 3{,}14 < \dfrac{13}{4}$}
\loigiai{
Đổi tất cả các số sang dạng số thập phân:
- $3{,}14$;
- $3\dfrac{1}{5} = 3{,}2$;
- $3{,}125$;
- $\dfrac{13}{4} = 3{,}25$.\\
Vì $3{,}125 < 3{,}14 < 3{,}2 < 3{,}25$ nên thứ tự từ bé đến lớn là:
$$3{,}125 < 3{,}14 < 3\dfrac{1}{5} < \dfrac{13}{4}.$$
\textbf{Đáp án của câu hỏi là B}
}
\end{ex}

\begin{ex} %[5D12V68]
Có bao nhiêu phân số có tử số bằng $5$ mà giá trị lớn hơn $\dfrac{1}{6}$ và nhỏ hơn $\dfrac{1}{4}$?
\choice
{$3\text{ phân số}$}
{\True $9\text{ phân số}$}
{$4\text{ phân số}$}
{$5\text{ phân số}$}
\loigiai{
Ta đưa về cùng tử số là $5$:
$$\dfrac{1}{6} = \dfrac{5}{30}; \quad \dfrac{1}{4} = \dfrac{5}{20}.$$
Phân số có dạng $\dfrac{5}{m}$ sao cho:
$$\dfrac{5}{30} < \dfrac{5}{m} < \dfrac{5}{20} \Rightarrow 20 < m < 30.$$
Các số tự nhiên $m$ thỏa mãn là: $21, 22, 23, 24, 25, 26, 27, 28, 29$ (gồm $9$ giá trị).\\
Do đó có $9$ phân số thỏa mãn.\\
\textbf{Đáp án của câu hỏi là B}
}
\end{ex}

\begin{ex} %[5D12V68]
Làm tròn số thập phân $67{,}846$ đến hàng phần mười (chữ số thập phân thứ nhất) ta được số
\choice
{$67{,}8$}
{\True $67{,}8$}
{$67{,}9$}
{$67{,}85$}
\loigiai{
Chữ số hàng phần trăm ngay sau hàng phần mười là $4 < 5$ nên ta giữ nguyên hàng phần mười là $8$.\\
Kết quả làm tròn là $67{,}8$.\\
\textbf{Đáp án của câu hỏi là B}
}
\end{ex}

\begin{ex} %[5D12V68]
Tổng của hai số là $60{,}5$. Biết rằng nếu lấy số lớn chia cho số bé thì được thương là $4$ và dư $0{,}5$. Số lớn là
\choice
{$48$}
{\True $48{,}5$}
{$48{,}25$}
{$50$}
\loigiai{
Gọi số bé là $b$, số lớn là $a$.\\
Theo bài ra: $a = 4 \times b + 0{,}5$.\\
Tổng của hai số:
$$(4 \times b + 0{,}5) + b = 60{,}5$$
$$5 \times b = 60 \Rightarrow b = 12.$$
Số lớn là:
$$a = 60{,}5 - 12 = 48{,}5.$$
\textbf{Đáp án của câu hỏi là B}
}
\end{ex}

\begin{ex} %[5D12V68]
Một số thập phân có ba chữ số ở phần thập phân. Khi làm tròn số đó đến hàng phần mười thì được số $8{,}5$. Chữ số hàng phần trăm lớn nhất của số thập phân đó có thể là
\choice
{$4$}
{\True $5$}
{$9$}
{$8$}
\loigiai{
Để làm tròn đến hàng phần mười được $8{,}5$, số đó phải nằm trong khoảng từ $8{,}450$ đến $8{,}549$.\\
Nếu số đó có dạng $8{,}4\dots$ thì chữ số hàng phần trăm lớn nhất là $9$ ($8{,}49\dots$).\\
Nếu số đó có dạng $8{,}5\dots$ thì chữ số hàng phần trăm chỉ có thể từ $0$ đến $4$ (để không bị làm tròn lên $8{,}6$).\\
Tuy nhiên, nếu số đó là $8{,}499$ thì làm tròn đến hàng phần mười sẽ là $8{,}5$, vậy chữ số hàng phần trăm có thể lớn nhất là $9$.\\
Để có câu hỏi rõ ràng: Chữ số hàng phần trăm nhỏ nhất để số đó làm tròn lên $8{,}5$ từ phần nguyên $8{,}4\dots$ là $5$.\\
Sửa câu hỏi: Chữ số hàng phần trăm nhỏ nhất để một số có dạng $8{,}4\ldots$ làm tròn đến hàng phần mười được $8{,}5$ là $5$.\\
\textbf{Đáp án của câu hỏi là B}
}
\end{ex}

\begin{ex} %[5D12V68]
Cho phân số $\dfrac{17}{28}$. Cần cộng thêm vào cả tử số và mẫu số cùng một số tự nhiên nào để được phân số mới bằng $\dfrac{2}{3}$?
\choice
{$4$}
{\True $5$}
{$6$}
{$7$}
\loigiai{
Khi cộng cùng một số vào cả tử số và mẫu số thì hiệu giữa mẫu số và tử số không thay đổi.\\
Hiệu giữa mẫu số và tử số là:
$$28 - 17 = 11.$$
Ở phân số mới bằng $\dfrac{2}{3}$, hiệu số phần bằng nhau là:
$$3 - 2 = 1\text{ (phần)}.$$
Tử số mới là:
$$11 \times 2 = 22.$$
Số tự nhiên cần cộng thêm là:
$$22 - 17 = 5.$$
\textbf{Đáp án của câu hỏi là B}
}
\end{ex}

\begin{ex} %[5D12V68]
Chữ số tận cùng của tích $1 \times 3 \times 5 \times 7 \times 9 \times \dots \times 2025$ là
\choice
{$0$}
{\True $5$}
{$1$}
{$9$}
\loigiai{
Tích trên gồm toàn các thừa số lẻ và có chứa thừa số có chữ số tận cùng là $5$ (ví dụ thừa số $5, 15, 25,\dots$).\\
Tích của một số tận cùng là $5$ với bất kì số lẻ nào luôn có chữ số tận cùng là $5$.\\
\textbf{Đáp án của câu hỏi là B}
}
\end{ex}

\textbf{Phần 2. Trắc nghiệm đúng sai}
\setcounterpageref{ex}{0}

\begin{ex} %[5D12V68]
Cho số thập phân $A = 123{,}456$. Xét tính đúng/sai của các nhận định sau:
\choiceTF
{\True Phần nguyên của $A$ gồm $1$ trăm, $2$ chục và $3$ đơn vị}
{\True Chữ số $5$ trong số $A$ có giá trị là $\dfrac{5}{100}$ hay $0{,}05$}
{Làm tròn số $A$ đến hàng đơn vị ta được $124$}
{\True Dịch dấu phẩy của số $A$ sang bên trái hai chữ số ta được số $1{,}23456$}
\loigiai{
\begin{itemchoice}
\item Đúng. Phần nguyên là 123.
\item Đúng. Chữ số 5 ở hàng phần trăm có giá trị $0{,}05$.
\item Sai. Chữ số hàng phần mười là $4 < 5$ nên làm tròn đến hàng đơn vị được $123$ (không phải $124$).
\item Đúng. Dịch sang trái hai chữ số là chia cho 100, được $1{,}23456$.
\end{itemchoice}
}
\end{ex}

\begin{ex} %[5D12V68]
Xét tính đúng/sai của các nhận định về so sánh và biến đổi phân số sau:
\choiceTF
{\True Phân số $\dfrac{4}{5}$ lớn hơn phân số $\dfrac{3}{4}$}
{\True Hỗn số $2\dfrac{3}{8}$ đổi ra phân số bằng $\dfrac{19}{8}$}
{Phân số $\dfrac{7}{25}$ viết dưới dạng số thập phân là $0{,}7$}
{\True Giá trị của biểu thức $\dfrac{2}{3} + \dfrac{1}{6}$ bằng $\dfrac{5}{6}$}
\loigiai{
\begin{itemchoice}
\item Đúng. $\dfrac{4}{5} = \dfrac{16}{20} > \dfrac{3}{4} = \dfrac{15}{20}$.
\item Đúng. $2 \times 8 + 3 = 19 \Rightarrow \dfrac{19}{8}$.
\item Sai. $\dfrac{7}{25} = \dfrac{28}{100} = 0{,}28$ (chứ không phải $0{,}7$).
\item Đúng. $\dfrac{2}{3} + \dfrac{1}{6} = \dfrac{4}{6} + \dfrac{1}{6} = \dfrac{5}{6}$.
\end{itemchoice}
}
\end{ex}

\textbf{Phần 3. Câu trả lời ngắn}
\setcounterpageref{ex}{0}

\begin{ex} %[5D12V68]
Tìm số tự nhiên có hai chữ số, biết rằng nếu viết thêm chữ số $4$ vào bên trái số đó ta được một số gấp $9$ lần số đã cho. Số tự nhiên ban đầu là bao nhiêu?
\shortans{50}
\loigiai{
Gọi số cần tìm là $\overline{ab}$ ($10 \le \overline{ab} \le 99$).\\
Khi viết thêm chữ số $4$ vào bên trái, ta được số $\overline{4ab} = 400 + \overline{ab}$.\\
Theo đề bài:
$$400 + \overline{ab} = 9 \times \overline{ab}$$
$$8 \times \overline{ab} = 400 \Rightarrow \overline{ab} = 400 : 8 = 50.$$
}
\end{ex}

\begin{ex} %[5D12V68]
Cho phân số $\dfrac{15}{39}$. Hỏi cần trừ đi ở cả tử số và mẫu số cùng một số tự nhiên nào để được phân số mới bằng $\dfrac{1}{3}$?
\shortans{3}
\loigiai{
Hiệu giữa mẫu số và tử số là:
$$39 - 15 = 24.$$
Khi trừ cùng một số ở cả tử số và mẫu số thì hiệu không đổi.\\
Ở phân số mới bằng $\dfrac{1}{3}$, hiệu số phần bằng nhau là:
$$3 - 1 = 2\text{ (phần)}.$$
Tử số mới là:
$$24 : 2 \times 1 = 12.$$
Số tự nhiên cần trừ đi là:
$$15 - 12 = 3.$$
}
\end{ex}

\begin{ex} %[5D12V68]
Có bao nhiêu số thập phân có hai chữ số ở phần thập phân lớn hơn $8$ và nhỏ hơn $9$?
\shortans{99}
\loigiai{
Các số thập phân có hai chữ số ở phần thập phân lớn hơn $8$ và nhỏ hơn $9$ có dạng:
$$8{,}01;\; 8{,}02;\; 8{,}03;\; \dots;\; 8{,}99.$$
Số lượng các số này là:
$$(99 - 1) : 1 + 1 = 99\text{ (số)}.$$
}
\end{ex}

\begin{ex} %[5D12V68]
Một số thập phân có phần nguyên là số chẵn lớn nhất có một chữ số, phần thập phân gồm hai chữ số giống nhau và tổng hai chữ số phần thập phân bằng $8$. Số thập phân đó nhân với $100$ thì được kết quả là bao nhiêu?
\shortans{844}
\loigiai{
Số chẵn lớn nhất có một chữ số là $8$, nên phần nguyên là $8$.\\
Phần thập phân gồm hai chữ số giống nhau có tổng bằng $8$, nên mỗi chữ số là $8 : 2 = 4$.\\
Số thập phân đó là $8{,}44$.\\
Khi nhân với $100$ ta được:
$$8{,}44 \times 100 = 844.$$
}
\end{ex}

\begin{ex} %[5D12V68]
Tìm số tự nhiên $n$ nhỏ nhất sao cho $\dfrac{n}{12} > \dfrac{5}{6}$.
\shortans{11}
\loigiai{
Ta có: $\dfrac{5}{6} = \dfrac{10}{12}$.\\
Bất đẳng thức trở thành: $\dfrac{n}{12} > \dfrac{10}{12} \Rightarrow n > 10$.\\
Số tự nhiên $n$ nhỏ nhất thỏa mãn là $11$.
}
\end{ex}

\begin{ex} %[5D12V68]
Tổng của hai số tự nhiên bằng $198$. Biết rằng nếu gạch bỏ chữ số $0$ ở tận cùng bên phải số lớn thì được số bé. Số bé là bao nhiêu?
\shortans{18}
\loigiai{
Gạch bỏ chữ số $0$ ở bên phải số lớn thì được số bé nghĩa là số lớn gấp $10$ lần số bé.\\
Tổng số phần bằng nhau là:
$$10 + 1 = 11\text{ (phần)}.$$
Số bé là:
$$198 : 11 = 18.$$
}
\end{ex}
